\chapter{Études de cas et dimensionnement}

% ============================================================
\section{Objectifs pédagogiques}
% ============================================================

À l'issue de ce chapitre, l'étudiant sera capable de :
\begin{itemize}
  \item Conduire une démarche complète de dimensionnement CVC en partant du cahier des charges ;
  \item Calculer les déperditions thermiques d'un bâtiment selon la méthode réglementaire (RE2020 / NF EN 12831) ;
  \item Sélectionner et dimensionner les équipements de production, distribution et émission de chaleur ;
  \item Dimensionner un réseau aéraulique et une centrale de traitement d'air (CTA) ;
  \item Appliquer les critères d'efficacité énergétique et vérifier la conformité réglementaire ;
  \item Intégrer la régulation et la GTB dans une étude globale ;
  \item Rédiger une note de calcul justifiée à présenter en situation professionnelle.
\end{itemize}

\textbf{Compétences mobilisées :} C1, C2, C3, C5, C7, C8, C9 \quad
\textbf{Savoirs associés :} S5.1, S8-A2, S8-A3, S8-A4, S8-A5, S8-B1, S8-B2, S8-B4, S8-B11

% ============================================================
\section{Introduction}
% ============================================================

Le technicien supérieur GCF est régulièrement amené à produire des études techniques complètes : appels d'offres, avant-projets sommaires (APS), avant-projets définitifs (APD), dossiers de consultation des entreprises (DCE). Ces études intègrent plusieurs disciplines enseignées séparément dans les chapitres précédents.

Ce chapitre propose trois études de cas représentatives :
\begin{enumerate}
  \item Un \textbf{immeuble résidentiel RE2020} — déperditions, production chaleur par PAC, distribution hydraulique, régulation ;
  \item Un \textbf{bâtiment tertiaire climatisé} — CTA double flux, réseau aéraulique, production froid, GTB ;
  \item Une \textbf{installation industrielle} — process thermique, sécurités, efficacité énergétique.
\end{enumerate}

\begin{definition}[title={Note de calcul}]
  La note de calcul est le document technique central d'une étude CVC. Elle regroupe :
  les données d'entrée (géométrie, occupation, conditions climatiques), les hypothèses
  retenues, les formules appliquées avec leurs références normatives, les résultats
  numériques et la sélection du matériel. Elle doit être \textbf{traçable, vérifiable et justifiée}.
\end{definition}

% ============================================================
\section{Méthodologie de dimensionnement CVC}
% ============================================================

\subsection{Démarche générale}

\begin{methode}[title={Étapes du dimensionnement CVC}]
  \begin{enumerate}
    \item \textbf{Collecte des données} : plans architectural, orientation, localisation climatique, destination des locaux, réglementation applicable.
    \item \textbf{Bilan thermique d'hiver} : calcul des déperditions par transmission et renouvellement d'air (NF EN 12831).
    \item \textbf{Bilan thermique d'été} : calcul des charges de refroidissement (apports solaires, internes, renouvellement).
    \item \textbf{Sélection des générateurs} : chaudière, PAC, groupe froid — selon besoins de pointe et modulabilité.
    \item \textbf{Distribution} : réseau hydraulique (dimensionnement des tuyauteries, pompes, équilibrage) et/ou aéraulique (gaines, ventilateurs, diffuseurs).
    \item \textbf{Émission} : radiateurs, planchers chauffants, ventilo-convecteurs, soufflage.
    \item \textbf{Régulation et GTB} : stratégies de commande, sondes, actionneurs.
    \item \textbf{Vérification RE2020} : calcul Bbio, Cep, DH.
    \item \textbf{Rédaction de la note de calcul} et sélection du matériel sur catalogue.
  \end{enumerate}
\end{methode}

\subsection{Données climatiques de référence}

Les températures extérieures de base $T_{ext,base}$ sont définies par zone climatique (NF EN 12831 et données Météo-France) :

\begin{center}
\begin{tabular}{llcc}
  \toprule
  Zone & Exemples de villes & $T_{ext,base}$ (\si{\celsius}) & DJU (base 18) \\
  \midrule
  H1a & Paris, Lille         & $-7$  & 2\,400 \\
  H1b & Reims, Nancy         & $-9$  & 2\,800 \\
  H1c & Strasbourg           & $-11$ & 3\,100 \\
  H2a & Nantes, La Rochelle  & $-4$  & 1\,800 \\
  H2b & Lyon, Grenoble       & $-8$  & 2\,400 \\
  H2c & Bordeaux, Toulouse   & $-3$  & 1\,600 \\
  H2d & Nice, Marseille      & 0     & 1\,200 \\
  H3   & Guadeloupe, Réunion  & +18  & 0 \\
  \bottomrule
\end{tabular}
\end{center}

\subsection{Rappel des formules fondamentales}

\subsubsection{Déperditions par transmission}

\begin{equation}
  \Phi_{T} = H_T \cdot (T_{int} - T_{ext,base})
\end{equation}

avec $H_T = \sum_{i} U_i \cdot A_i$ le coefficient de déperdition global par transmission (\si{\watt\per\kelvin}).

\subsubsection{Déperditions par ventilation et infiltrations}

\begin{equation}
  \Phi_{V} = H_V \cdot (T_{int} - T_{ext,base})
\end{equation}

avec $H_V = \rho_{air} \cdot c_{p,air} \cdot q_v$ (\si{\watt\per\kelvin}), où $q_v$ est le débit volumique d'air neuf (\si{\cubic\meter\per\second}).

\subsubsection{Puissance de chauffage de pointe}

\begin{equation}
  \Phi_{ch} = \Phi_{T} + \Phi_{V} \quad [\si{\watt}]
\end{equation}

\subsubsection{Coefficient de performance de la PAC}

\begin{equation}
  \mathrm{COP} = \frac{P_{calorifique}}{P_{electrique}}
\end{equation}

\subsubsection{Perte de charge en réseau hydraulique}

\begin{equation}
  \Delta P = \lambda \cdot \frac{L}{D} \cdot \frac{\rho v^2}{2} + \sum \xi_i \cdot \frac{\rho v^2}{2} \quad [\si{\pascal}]
\end{equation}

\subsubsection{Puissance frigorifique nécessaire}

\begin{equation}
  Q_f = Q_{apports\_solaires} + Q_{apports\_internes} + Q_{ventilation}
\end{equation}

% ============================================================
\section{Étude de cas 1 — Immeuble résidentiel RE2020}
% ============================================================

\subsection{Présentation du projet}

\begin{exemple}[title={Contexte}]
  Immeuble résidentiel collectif R+3, zone climatique H1a (Paris), altitude \SI{75}{\metre}.
  Surface habitable totale : \SI{1200}{\metre\squared} (12 appartements de \SI{100}{\metre\squared}).
  Destination : logements sociaux BBC-RE2020.
  Système CVC retenu : PAC air/eau collective + planchers chauffants + VMC double flux.
\end{exemple}

\subsection{Données architecturales}

\begin{center}
\begin{tabular}{llccc}
  \toprule
  Paroi & Description & Surface (\si{\metre\squared}) & $U$ (\si{\watt\per\metre\squared\per\kelvin}) & $U \cdot A$ (\si{\watt\per\kelvin}) \\
  \midrule
  Murs ext.    & Béton + isolation ITE 14 cm & 620 & 0,22 & 136,4 \\
  Toiture      & Terrasse accessible isolée   & 300 & 0,15 & 45,0  \\
  Plancher bas & Dalle sur terre-plein        & 300 & 0,30 & 90,0  \\
  Vitrages     & Double vitrage Uw 1,2        & 180 & 1,20 & 216,0 \\
  Portes ext.  & Porte palière isolée         & 30  & 1,40 & 42,0  \\
  \midrule
  \textbf{Total} & & & & \textbf{529,4} \\
  \bottomrule
\end{tabular}
\end{center}

\subsection{Calcul des déperditions}

\textbf{Données :} $T_{int} = \SI{19}{\celsius}$, $T_{ext,base} = \SI{-7}{\celsius}$, $\Delta T = \SI{26}{\kelvin}$.

\textbf{Renouvellement d'air :}
\begin{itemize}
  \item Débit VMC double flux : $q_v = \SI{0,5}{\per\hour} \times \SI{3600}{\metre\cubed} = \SI{1800}{\metre\cubed\per\hour}$
  \item Rendement échangeur : $\eta = 0,85$ (75\,\% récupéré)
  \item $H_V = 0,34 \times 1800 \times (1 - 0,85) = \SI{91,8}{\watt\per\kelvin}$
\end{itemize}

\textbf{Coefficient de déperdition total :}
\begin{equation}
  H_{total} = H_T + H_V = 529{,}4 + 91{,}8 = \SI{621{,}2}{\watt\per\kelvin}
\end{equation}

\textbf{Puissance de chauffage de pointe :}
\begin{equation}
  \Phi_{ch} = 621{,}2 \times 26 = \SI{16{,}15}{\kilo\watt}
\end{equation}

\begin{attention}
  La puissance de pointe calculée correspond aux déperditions en conditions hivernales extrêmes.
  Elle ne prend pas en compte les apports solaires et internes, qui sont favorables en hiver
  mais doivent être vérifiés en été pour le confort d'été (calcul des charges de refroidissement).
\end{attention}

\subsection{Sélection de la PAC}

\begin{methode}[title={Critères de sélection d'une PAC collective}]
  \begin{enumerate}
    \item Puissance calorifique $\geq \Phi_{ch}$ (avec marge de 15\,\% pour pertes réseau)
    \item COP $\geq 3{,}2$ à $-7/35$ (SCOP annuel selon EN 14825)
    \item Fluide frigorigène à faible GWP ($<$ 750 — exigences F-gas 2024)
    \item Compatibilité avec plancher chauffant (eau à \SI{35}{\celsius} max)
    \item Niveau acoustique $\leq$ 50 dB(A) à \SI{1}{\metre} (mitoyenneté)
  \end{enumerate}
\end{methode}

\textbf{Puissance installée avec marge :}
\begin{equation}
  P_{PAC} = 16{,}15 \times 1{,}15 = \SI{18{,}6}{\kilo\watt} \rightarrow \text{sélection : PAC 20 kW}
\end{equation}

\textbf{Consommation électrique annuelle estimée :}
\begin{equation}
  W_{el} = \frac{P_{calorifique} \times DJU \times 24}{SCOP \times \Delta T_{base}} = \frac{20 \times 2400 \times 24}{3{,}5 \times 26} \approx \SI{12\,615}{\kilo\watt\hour\per\year}
\end{equation}

\subsection{Dimensionnement du plancher chauffant}

Le plancher chauffant basse température (PCBT) délivre la chaleur par rayonnement depuis la dalle.

\textbf{Puissance surfacique utile :}
\begin{equation}
  q_{plancher} = \frac{\Phi_{ch}}{A_{totale}} = \frac{16\,150}{1200} = \SI{13{,}5}{\watt\per\metre\squared}
\end{equation}

\textbf{Température de départ eau :} \SI{35}{\celsius} (retour \SI{28}{\celsius}), compatible PAC.

\textbf{Pas de pose des tubes :} espacement $e = \SI{15}{\centi\metre}$ (standard pour habitat). Débit par circuit :
\begin{equation}
  \dot{m}_{circuit} = \frac{q_{plancher} \times A_{circuit}}{\rho \cdot c_p \cdot \Delta T_{eau}} = \frac{13{,}5 \times 15}{1000 \times 4186 \times 7 / 3600} \approx \SI{37{,}4}{\litre\per\hour}
\end{equation}

\subsection{Réseau hydraulique}

\begin{center}
\begin{tabular}{lcccc}
  \toprule
  Tronçon & Débit (\si{\metre\cubed\per\hour}) & $\Delta P_{lin}$ (\si{\pascal\per\metre}) & Longueur (\si{\metre}) & $\Delta P_{tot}$ (\si{\pascal}) \\
  \midrule
  Collecteur$\to$R+3 & 1,8 & 120 & 15 & 1\,800 \\
  R+3$\to$appart 1  & 0,15 & 80  & 12 & 960 \\
  Circuit PCBT       & 0,15 & 60  & 60 & 3\,600 \\
  \midrule
  \textbf{Circuit le plus défavorable} & & & & \textbf{6\,360} \\
  \bottomrule
\end{tabular}
\end{center}

\textbf{Sélection de la pompe :} $q_v = \SI{1,8}{\metre\cubed\per\hour}$, $\Delta P = \SI{6\,360}{\pascal} \approx \SI{6,4}{\metre} CE$.

\subsection{Vérification RE2020}

\begin{definition}[title={Indicateurs RE2020 à vérifier}]
  \begin{itemize}
    \item \textbf{Bbio} (Besoin bioclimatique) : efficacité de l'enveloppe — $\leq Bbio_{max}$
    \item \textbf{Cep} (Consommation d'énergie primaire) : $\leq Cep_{max}$
    \item \textbf{Cep,nr} (Énergie primaire non renouvelable) : $\leq Cep,nr_{max}$
    \item \textbf{Icénergie} (Impact carbone énergie) : $\leq$ seuil progressif
    \item \textbf{DH} (Degrés-Heures de confort d'été) : $\leq 1\,250$ DH (résidentiel)
  \end{itemize}
\end{definition}

Pour notre immeuble PAC (coefficient de conversion électricité : $f_{prim} = 2{,}3$) :
\begin{equation}
  Cep = \frac{W_{el} \times f_{prim}}{A_{ref}} = \frac{12\,615 \times 2{,}3}{1\,200} = \SI{24{,}2}{\kilo\watt\hour\per\metre\squared\per\year}
\end{equation}

Ce résultat est très inférieur au seuil RE2020 de \SI{70}{\kilo\watt\hour_{ep}\per\metre\squared\per\year} pour le chauffage en zone H1a.

% ============================================================
\section{Étude de cas 2 — Bâtiment tertiaire climatisé}
% ============================================================

\subsection{Présentation du projet}

\begin{exemple}[title={Contexte}]
  Immeuble de bureaux R+5, zone H2b (Lyon), surface utile : \SI{3\,000}{\metre\squared}.
  Occupation : 300 personnes, bureaux paysagers et salles de réunion.
  Système CVC : groupe d'eau glacée (GEG) + CTA double flux avec batterie froide
  + ventilo-convecteurs (VRV en complément), GTB centralisée.
  Objectif : confort d'été garanti ($T_{int} \leq \SI{26}{\celsius}$), labelisation HQE.
\end{exemple}

\subsection{Calcul des charges de refroidissement}

\textbf{Apports internes (été) :}
\begin{center}
\begin{tabular}{lcc}
  \toprule
  Source & Valeur unitaire & Total \\
  \midrule
  Personnes (sensible) & $70\,\si{\watt}/pers$ & $300 \times 70 = \SI{21\,000}{\watt}$ \\
  Éclairage LED & $8\,\si{\watt\per\metre\squared}$ & $3000 \times 8 = \SI{24\,000}{\watt}$ \\
  Bureautique & $150\,\si{\watt}/poste$ & $200 \times 150 = \SI{30\,000}{\watt}$ \\
  \midrule
  \textbf{Total apports internes} & & \SI{75\,000}{\watt} \\
  \bottomrule
\end{tabular}
\end{center}

\textbf{Apports solaires (façade Sud, fenêtres \SI{600}{\metre\squared}, $g = 0{,}35$, $G_{max} = \SI{600}{\watt\per\metre\squared}$) :}
\begin{equation}
  Q_{solaire} = g \times A_{vitrage} \times G_{max} = 0{,}35 \times 600 \times 600 = \SI{126\,000}{\watt}
\end{equation}

\textbf{Apports par renouvellement d'air (conditions : $T_{ext} = \SI{33}{\celsius}$, $T_{int} = \SI{26}{\celsius}$) :}
\begin{equation}
  q_v = 300\,pers \times \SI{25}{\litre\per\second\per\personne} = \SI{7{,}5}{\metre\cubed\per\second}
\end{equation}
\begin{equation}
  Q_{vent} = \rho \cdot c_p \cdot q_v \cdot (T_{ext} - T_{int}) = 1{,}2 \times 1005 \times 7{,}5 \times 7 = \SI{63\,315}{\watt}
\end{equation}

\textbf{Charge totale de refroidissement :}
\begin{equation}
  Q_f = 75\,000 + 126\,000 + 63\,315 = \SI{264\,315}{\watt} \approx \SI{265}{\kilo\watt}
\end{equation}

\subsection{Sélection du groupe d'eau glacée}

\textbf{Puissance frigorifique nécessaire (marge 10\,\%) :}
\begin{equation}
  Q_{GEG} = 265 \times 1{,}10 = \SI{292}{\kilo\watt} \rightarrow \text{GEG 300 kW}
\end{equation}

\textbf{COP du groupe froid à $5/12$\,\si{\celsius} :} $\mathrm{EER} = 3{,}0$ (classe A, norme EN 14511).

\textbf{Puissance absorbée :}
\begin{equation}
  P_{el} = \frac{Q_{GEG}}{EER} = \frac{300}{3{,}0} = \SI{100}{\kilo\watt}
\end{equation}

\textbf{Circuit eau glacée :} température $5/12\,\si{\celsius}$, débit :
\begin{equation}
  \dot{m}_{EG} = \frac{Q_{GEG}}{\rho \cdot c_p \cdot \Delta T} = \frac{300\,000}{1000 \times 4186 \times 7} = \SI{10{,}2}{\metre\cubed\per\hour}
\end{equation}

\subsection{Dimensionnement de la CTA double flux}

\begin{methode}[title={Dimensionnement d'une CTA}]
  \begin{enumerate}
    \item Déterminer le débit d'air neuf $q_{v,AN}$ (réglementation aéraulique, besoins hygiéniques).
    \item Choisir les températures de soufflage ($T_{soufflage}$) et de reprise.
    \item Dimensionner la batterie froide : $Q_{batterie} = \rho \cdot c_p \cdot q_v \cdot (T_{reprise} - T_{soufflage})$.
    \item Sélectionner le ventilateur (débit + pression disponible).
    \item Vérifier le bilan enthalpique sur le diagramme de l'air humide.
  \end{enumerate}
\end{methode}

\textbf{Débit air total CTA :} $q_{v,total} = \SI{25\,000}{\metre\cubed\per\hour}$ (8 vol/h + hygiénique).

\textbf{Température de soufflage : \SI{15}{\celsius}} (reprise \SI{26}{\celsius}).

\textbf{Puissance de la batterie froide de la CTA :}
\begin{equation}
  Q_{batterie} = \frac{1{,}2 \times 1005 \times 25\,000}{3600} \times (26 - 15) = \SI{92\,125}{\watt} \approx \SI{92}{\kilo\watt}
\end{equation}

\textbf{Perte de charge réseau aéraulique (circuit le plus défavorable) :}
\begin{center}
\begin{tabular}{lccc}
  \toprule
  Tronçon & Débit (\si{\metre\cubed\per\hour}) & $\Delta P$ (\si{\pascal}) & Vitesse (\si{\metre\per\second}) \\
  \midrule
  Gaine principale & 25\,000 & 180 & 6,0 \\
  Gaines secondaires & 5\,000 & 95 & 4,5 \\
  Diffuseurs & — & 20 & — \\
  CTA (filtres+batteries) & — & 350 & — \\
  \midrule
  \textbf{Total} & & \textbf{645} & \\
  \bottomrule
\end{tabular}
\end{center}

\textbf{Sélection du ventilateur :} $q_v = \SI{25\,000}{\metre\cubed\per\hour}$, $\Delta P_{tot} = \SI{645}{\pascal}$.
Puissance ventilateur (rendement $\eta = 0{,}75$) :
\begin{equation}
  P_{vent} = \frac{q_v \times \Delta P_{tot}}{\eta} = \frac{(25\,000/3600) \times 645}{0{,}75} = \SI{5\,972}{\watt} \approx \SI{6}{\kilo\watt}
\end{equation}

\subsection{Intégration GTB}

\begin{definition}[title={Architecture GTB pour bâtiment tertiaire}]
  La GTB centralise la supervision et la commande de toutes les installations CVC :
  \begin{itemize}
    \item \textbf{Niveau terrain} : capteurs (sondes de température, CO\textsubscript{2}, humidité, pression), actionneurs (vannes, volets, variateurs de vitesse) ;
    \item \textbf{Niveau automation} : automates DDC (Direct Digital Control) gérant les boucles de régulation locales ;
    \item \textbf{Niveau supervision} : IHM de surveillance et de pilotage, archivage des données, alarmes.
  \end{itemize}
  Protocoles courants : BACnet (IP ou MS/TP), Modbus RTU/TCP, KNX.
\end{definition}

\textbf{Fonctions GTB pour ce projet :}
\begin{itemize}
  \item Programmation horaire des CTA (démarrage anticipé, relance matinale) ;
  \item Régulation de la température d'eau glacée en fonction de la charge (reset) ;
  \item Free-cooling aéraulique ($T_{ext} < T_{int} - \SI{2}{\celsius}$) ;
  \item Comptage des consommations électriques (GEG, CTA, éclairage) ;
  \item Alarmes de dépassement de température et de pression.
\end{itemize}

% ============================================================
\section{Étude de cas 3 — Installation industrielle}
% ============================================================

\subsection{Présentation du projet}

\begin{exemple}[title={Contexte}]
  Atelier de production pharmaceutique, surface : \SI{800}{\metre\squared}.
  Exigences : température $22 \pm \SI{1}{\celsius}$, humidité relative $50 \pm 5\,\%$,
  surpression par rapport aux zones adjacentes, renouvellement d'air conforme BPF
  (Bonnes Pratiques de Fabrication — 20 volumes/heure minimum).
  Énergie : vapeur process à \SI{6}{\bar} disponible sur site, eau glacée \SI{6/12}{\celsius}.
\end{exemple}

\subsection{Bilan thermique et massique}

\textbf{Charges internes spécifiques (process) :}
\begin{center}
\begin{tabular}{lcc}
  \toprule
  Équipement & Puissance dissipée (\si{\kilo\watt}) & Fonctionnement \\
  \midrule
  Malaxeurs (4 $\times$ 15 kW) & 45 (rendement 75\,\%) & 16\,h/j \\
  Éclairage salles blanches    & 12 & 16\,h/j \\
  Personnel (50 pers $\times$ 100 W) & 5 & 8\,h/j \\
  Pertes parois (isolation $U_{moy}=0,3$) & 8 & continu \\
  \midrule
  \textbf{Total pointe} & \textbf{70} & \\
  \bottomrule
\end{tabular}
\end{center}

\textbf{Débit d'air neuf (20 vol/h) :}
\begin{equation}
  q_{v,AN} = 20 \times \SI{800}{\metre\squared} \times \SI{3}{\metre} = \SI{48\,000}{\metre\cubed\per\hour}
\end{equation}

\textbf{Charge frigorifique totale (été, $T_{ext} = \SI{33}{\celsius}$) :}
\begin{equation}
  Q_{f,total} = Q_{process} + Q_{vent} = 70 + \frac{1{,}2 \times 1005 \times (48\,000/3600) \times (33-22)}{1000} = 70 + \SI{176}{\kilo\watt} = \SI{246}{\kilo\watt}
\end{equation}

\subsection{Traitement d'air et hygrométrie}

\begin{methode}[title={Contrôle hygrométrique en salle propre}]
  \begin{enumerate}
    \item Calculer la teneur en vapeur d'eau requise : $w_{cible} = f(T_{int}, HR_{cible})$ — diagramme de l'air humide.
    \item Déshumidification si $w_{ext} > w_{cible}$ : batterie froide amenée à $T_{rosée}$, puis remontée en température.
    \item Humidification si $w_{ext} < w_{cible}$ : humidificateur vapeur (vapeur propre, sans additifs chimiques en pharmaceutique).
    \item Régulation proportionnelle-intégrale de l'humidité relative avec sonde HR (\SI{\pm 1}{\percent} de précision).
  \end{enumerate}
\end{methode}

\textbf{Calcul de la teneur en eau à $22\,\si{\celsius}$, $HR = 50\,\%$ :}
\begin{equation}
  P_{vs}(22°C) = \SI{2\,643}{\pascal}
\end{equation}
\begin{equation}
  w = 0{,}622 \times \frac{HR \times P_{vs}}{P_{atm} - HR \times P_{vs}} = 0{,}622 \times \frac{0{,}5 \times 2643}{101\,325 - 0{,}5 \times 2643} = \SI{8{,}2}{\gram\per\kilogram}
\end{equation}

\subsection{Efficacité énergétique et économies d'énergie}

\begin{definition}[title={PUE et indicateurs performance installation industrielle}]
  \begin{itemize}
    \item \textbf{SEER} (Seasonal Energy Efficiency Ratio) : EER annuel moyenné sur les conditions climatiques.
    \item \textbf{Récupération de chaleur} sur air extrait : réducteur des consommations de chauffage/humidification.
    \item \textbf{Variateurs de vitesse} sur ventilateurs et pompes : économie proportionnelle au cube de la réduction de vitesse.
    \item \textbf{Free-cooling} : utilisation de l'air extérieur frais pour le refroidissement sans compresseur.
  \end{itemize}
\end{definition}

\textbf{Économie par récupération de chaleur sur air extrait (échangeur à plaques, $\eta = 0,75$) :}
\begin{equation}
  Q_{recup} = 0{,}75 \times \frac{1{,}2 \times 1005 \times 48\,000}{3600} \times (22 - (-5)) = \SI{362}{\kilo\watt}
\end{equation}

Cette économie permet de réduire considérablement la puissance de la batterie chaude vapeur en hiver.

\textbf{Économie annuelle estimée (chauffage et humidification) :}
\begin{equation}
  W_{econo} = 362 \times 3\,000\,h/an \times 10^{-3} \approx \SI{1\,086}{\mega\watt\hour\per\year}
\end{equation}

\subsection{Sécurités spécifiques installation industrielle}

\begin{attention}
  En milieu industriel et pharmaceutique, les sécurités CVC sont critiques pour la
  production et la sécurité des personnels :
  \begin{itemize}
    \item \textbf{Surpression salle blanche} : maintenir $+15\,\si{\pascal}$ par rapport aux zones adjacentes (évite contamination).
    \item \textbf{Alarme température} : dépassement de $22 \pm \SI{2}{\celsius}$ déclenche une alarme de niveau 2 avec traçabilité.
    \item \textbf{Filtre à haute efficacité} : filtration HEPA H14 ($\geq 99{,}995\,\%$ de rétention à \SI{0,3}{\micro\metre}).
    \item \textbf{Redondance des équipements} : doublement des ventilateurs (N+1), groupe froid de secours.
    \item \textbf{Qualification IQ/OQ/PQ} : installation CVC qualifiée selon les référentiels BPF (FDA, EMA).
  \end{itemize}
\end{attention}

% ============================================================
\section{Éléments de réglementation}
% ============================================================

\subsection{Réglementation thermique RE2020}

La RE2020 (entrée en vigueur le 1\textsuperscript{er} janvier 2022) fixe des exigences de performance énergétique et environnementale pour les bâtiments neufs :

\begin{center}
\begin{tabular}{lp{7cm}l}
  \toprule
  Indicateur & Définition & Exigence type \\
  \midrule
  Bbio & Besoin bioclimatique (enveloppe seule) & $\leq Bbio_{max}$ \\
  Cep & Consommation énergie primaire totale & $\leq Cep_{max}$ \\
  Cep,nr & Part non renouvelable de Cep & $\leq Cep,nr_{max}$ \\
  Icénergie & Impact CO\textsubscript{2} des consommations & $\leq$ seuil progressif \\
  Iconstruction & Impact CO\textsubscript{2} des matériaux de construction & $\leq$ seuil progressif \\
  DH & Degrés-Heures inconfort d'été & $\leq 1\,250$ DH résidentiel \\
  \bottomrule
\end{tabular}
\end{center}

\subsection{Réglementation aéraulique et qualité de l'air intérieur (QAI)}

\begin{itemize}
  \item \textbf{Arrêté du 24 mars 1982} (modifié) : débits minimaux d'air neuf par local et par usage.
  \item \textbf{Décret n°2012-14} : obligation de vérification des systèmes de ventilation mécanique.
  \item \textbf{NF EN 16798-1} : paramètres d'entrée pour la conception des bâtiments (QAI, catégories IDA 1 à 4).
  \item \textbf{Arrêté 2011 Établissements recevant du public} : surveillance CO\textsubscript{2} et COVT obligatoire dans les ERP de catégories 1 à 4.
\end{itemize}

\subsection{Fluides frigorigènes et réglementation F-gas}

\begin{definition}[title={Règlement F-gas (UE) n°517/2014 et révision 2024}]
  Calendrier de réduction progressive des HFC à fort PRG (Potentiel de Réchauffement Global) :
  \begin{itemize}
    \item \textbf{2025} : interdiction équipements neufs charge $>$ 3 kg avec GWP $>$ 750.
    \item \textbf{2027} : interdiction GWP $>$ 150 pour équipements hermétiques.
    \item \textbf{2032} : interdiction générale HFC pour équipements neufs dans de nombreuses applications.
  \end{itemize}
  Alternatives : R-290 (propane, GWP=3), R-32 (GWP=675), R-1234yf (GWP<1), R-744 (CO\textsubscript{2}, GWP=1).
\end{definition}

\begin{attention}
  Toute intervention sur un circuit frigorifique de charge $\geq \SI{5}{\tonne}$ CO\textsubscript{2}
  équivalent nécessite une attestation de capacité délivrée par un organisme accrédité
  (Certifrigo, Qualifrigo). L'enregistrement des fluides récupérés est obligatoire.
\end{attention}

\subsection{Normes de conception}

\begin{center}
\begin{tabular}{ll}
  \toprule
  Norme & Domaine \\
  \midrule
  NF EN 12831-1 & Calcul des déperditions thermiques de base \\
  NF EN 12831-3 & Calcul des besoins de chauffage des locaux \\
  NF EN 15243 & Calcul de la charge de froid \\
  NF EN 14511 & Performances des climatiseurs et PAC \\
  NF EN 16798-3 & Conception des systèmes de ventilation \\
  DTU 65.11 & Dispositifs de sécurité installations chauffage \\
  DTU 65.12 & Réseaux de tuyauteries en acier \\
  \bottomrule
\end{tabular}
\end{center}
